\recipe{Bella Italia’s own Panzerotti Fungi}
\label{recipe:panzerotti_fungi}
\makes{2-3 portions}
\worktime{\SI{30} \minute}
\totaltime{\SI{40} \minute}
\source{Bella Italia}

(no guys, it’s actually really good)

\begin{ingredients}

  \item 1 pack fresh pasta sheets 
  \item 200g mushrooms, finely chopped
  \item 1 clove garlic, chopped
  \item 1 tbsp olive oil
  \item 100g ricotta
  \item 2 tbsp grated parmesan
  \item Salt and black pepper
  \item 40g butter
\end{ingredients}


\begin{method}
Put a pan on medium heat with the olive oil.

Add garlic, then mushrooms and a pinch of salt. Cook for 5–7
minutes until they look dry (not watery).

Let the mushrooms cool slightly, then mix with ricotta and
parmesan.Season with salt and pepper.

Lay out your cooked pasta sheets. Add small spoonfuls of filling,
fold over, and press the edges closed.

Can be scruffy - doesn’t have to be perfect!

Boil in salted water for 2–3 minutes.

Melt butter in a pan, add the cooked pasta, toss gently, serve.
  
\end{method}


\begin{figure}[h]
  	\centering
  	
  	\begin{tikzpicture}
  		% Save the image in a box to get its dimensions
  		\sbox0{\includegraphics[width = 0.6\linewidth]{images/MissingImage.png}}%
  
  		% Draw a clipping rectangle with rounded corners
  		\path[clip, rounded corners=0.5cm] (0,0) rectangle (\wd0,\ht0);
  
  		% Place the image inside the clipped area
  		\node[inner sep=0pt, outer sep=0pt, anchor=south west] at (0,0) {\usebox0};
	\end{tikzpicture}

  \caption*{caption for image with date and information, 1st January 2026}
\end{figure}

\begin{notes}
TIPS
If the mushrooms look dry, you’ve done it right
If a few burst open, it doesn’t matter
Butter + cheese will carry the flavour! 
Key to this recipe is to persistently remind your friends how good your meal is. Never give up the act.
\end{notes}

